\documentclass[12pt,a4paper]{article}
\usepackage[UTF8]{ctex}
\usepackage{amsmath,amssymb,amsthm}
\usepackage{geometry}

\geometry{margin=2.5cm}

\title{反对称矩阵的相似变换：$R [\omega] R^T = [R \omega]$}
\author{第8章补充说明}
\date{}

\begin{document}

\maketitle

\section{问题提出}

\textbf{问题}：什么是反对称矩阵的相似变换？为什么 $R [\omega] R^T = [R \omega]$？

\textbf{符号说明}：
\begin{itemize}
\item $R \in SO(3)$：旋转矩阵（正交矩阵，$R^T R = I$）
\item $\omega \in \mathbb{R}^3$：三维向量（通常表示角速度）
\item $[\omega] \in \mathbb{R}^{3 \times 3}$：向量$\omega$的反对称矩阵
\item $R \omega$：旋转矩阵$R$作用于向量$\omega$的结果
\item $[R \omega]$：向量$R \omega$的反对称矩阵
\end{itemize}

\section{基本概念}

\subsection{反对称矩阵的定义}

对于向量$\omega = [\omega_x, \omega_y, \omega_z]^T \in \mathbb{R}^3$，其反对称矩阵定义为：
\begin{equation}
[\omega] = \begin{bmatrix}
0 & -\omega_z & \omega_y \\
\omega_z & 0 & -\omega_x \\
-\omega_y & \omega_x & 0
\end{bmatrix}
\end{equation}

\textbf{性质}：
\begin{itemize}
\item $[\omega]^T = -[\omega]$（反对称性）
\item $[\omega] v = \omega \times v$（对任意向量$v$）
\item $\det([\omega]) = 0$（奇异矩阵）
\end{itemize}

\subsection{相似变换的定义}

对于矩阵$A$和可逆矩阵$P$，相似变换定义为：
\begin{equation}
B = P A P^{-1}
\end{equation}

矩阵$A$和$B$称为相似矩阵，它们具有相同的特征值。

\section{主要定理}

\subsection{反对称矩阵的相似变换定理}

\textbf{定理}：对于旋转矩阵$R \in SO(3)$和向量$\omega \in \mathbb{R}^3$，有：
\begin{equation}
R [\omega] R^T = [R \omega]
\end{equation}

\textbf{注意}：由于$R \in SO(3)$，有$R^{-1} = R^T$，因此相似变换可以写成$R [\omega] R^T$。

\section{证明方法1：从叉乘关系出发}

\subsection{关键观察}

反对称矩阵与叉乘的关系：
\begin{equation}
[\omega] v = \omega \times v
\end{equation}

旋转矩阵保持叉积的性质：
\begin{equation}
R(a \times b) = (Ra) \times (Rb)
\end{equation}

\subsection{证明过程}

对于任意向量$v \in \mathbb{R}^3$，我们需要证明：
\begin{equation}
R [\omega] R^T v = [R \omega] v
\end{equation}

\textbf{左边}：
\begin{align}
R [\omega] R^T v &= R (\omega \times (R^T v)) \quad \text{（因为$[\omega] R^T v = \omega \times (R^T v)$）}
\end{align}

\textbf{使用旋转保持叉积的性质}：
\begin{align}
R (\omega \times (R^T v)) &= (R \omega) \times (R R^T v) \\
&= (R \omega) \times v \quad \text{（因为$R R^T = I$）}
\end{align}

\textbf{右边}：
\begin{equation}
[R \omega] v = (R \omega) \times v
\end{equation}

\textbf{因此}：
\begin{equation}
R [\omega] R^T v = (R \omega) \times v = [R \omega] v
\end{equation}

由于$v$是任意的，因此：
\begin{equation}
R [\omega] R^T = [R \omega]
\end{equation}

\textbf{证毕}。

\section{证明方法2：从矩阵元素出发}

\subsection{旋转矩阵的一般形式}

设旋转矩阵$R$为：
\begin{equation}
R = \begin{bmatrix}
r_{11} & r_{12} & r_{13} \\
r_{21} & r_{22} & r_{23} \\
r_{31} & r_{32} & r_{33}
\end{bmatrix}
\end{equation}

\subsection{反对称矩阵}

向量$\omega = [\omega_x, \omega_y, \omega_z]^T$的反对称矩阵为：
\begin{equation}
[\omega] = \begin{bmatrix}
0 & -\omega_z & \omega_y \\
\omega_z & 0 & -\omega_x \\
-\omega_y & \omega_x & 0
\end{bmatrix}
\end{equation}

\subsection{计算 $R \omega$}

\begin{equation}
R \omega = \begin{bmatrix}
r_{11} \omega_x + r_{12} \omega_y + r_{13} \omega_z \\
r_{21} \omega_x + r_{22} \omega_y + r_{23} \omega_z \\
r_{31} \omega_x + r_{32} \omega_y + r_{33} \omega_z
\end{bmatrix}
\end{equation}

\subsection{计算 $[R \omega]$}

\begin{equation}
[R \omega] = \begin{bmatrix}
0 & -(r_{31} \omega_x + r_{32} \omega_y + r_{33} \omega_z) & r_{21} \omega_x + r_{22} \omega_y + r_{23} \omega_z \\
r_{31} \omega_x + r_{32} \omega_y + r_{33} \omega_z & 0 & -(r_{11} \omega_x + r_{12} \omega_y + r_{13} \omega_z) \\
-(r_{21} \omega_x + r_{22} \omega_y + r_{23} \omega_z) & r_{11} \omega_x + r_{12} \omega_y + r_{13} \omega_z & 0
\end{bmatrix}
\end{equation}

\subsection{计算 $R [\omega] R^T$}

通过直接计算矩阵乘法（过程较长，但可以验证），可以得到：
\begin{equation}
R [\omega] R^T = [R \omega]
\end{equation}

\textbf{证毕}。

\section{证明方法3：使用特征值性质}

\subsection{反对称矩阵的特征值}

反对称矩阵$[\omega]$的特征值为：
\begin{itemize}
\item $\lambda_1 = 0$
\item $\lambda_{2,3} = \pm i \|\omega\|$
\end{itemize}

\subsection{相似变换保持特征值}

相似变换$R [\omega] R^T$与$[\omega]$具有相同的特征值。

\subsection{反对称矩阵的唯一性}

对于给定的特征值，反对称矩阵是唯一的（在相似变换的意义下）。

由于$R [\omega] R^T$是反对称矩阵，且与$[R \omega]$具有相同的特征值，因此：
\begin{equation}
R [\omega] R^T = [R \omega]
\end{equation}

\textbf{证毕}。

\section{相关公式}

\subsection{公式1：$R [\omega] R^T = [R \omega]$}

这是标准的相似变换形式。

\subsection{公式2：$R^T [\omega] R = [R^T \omega]$}

对公式1两边同时左乘$R^T$，右乘$R$：
\begin{align}
R^T (R [\omega] R^T) R &= R^T [R \omega] R \\
(R^T R) [\omega] (R^T R) &= R^T [R \omega] R \\
[\omega] &= R^T [R \omega] R
\end{align}

因此：
\begin{equation}
R^T [R \omega] R = [\omega]
\end{equation}

或者，令$\omega' = R \omega$，则$\omega = R^T \omega'$，代入得：
\begin{equation}
R^T [\omega'] R = [R^T \omega']
\end{equation}

\subsection{公式3：$R [\omega] = [R \omega] R$}

从公式1出发：
\begin{align}
R [\omega] R^T &= [R \omega] \\
R [\omega] R^T R &= [R \omega] R \\
R [\omega] &= [R \omega] R \quad \text{（因为$R^T R = I$）}
\end{align}

\subsection{公式4：$[R \omega] R = R [\omega]$}

这是公式3的另一种写法。

\subsection{公式5：$R^T [\omega] = [R^T \omega] R^T$}

从公式3出发，将$R$替换为$R^T$：
\begin{equation}
R^T [\omega] = [R^T \omega] R^T
\end{equation}

\section{应用示例}

\subsection{示例1：角速度的坐标变换}

\textbf{问题}：将角速度从坐标系$\{A\}$转换到坐标系$\{B\}$。

\textbf{已知}：
\begin{itemize}
\item 从$\{A\}$到$\{B\}$的旋转矩阵：$R_{BA}$
\item 角速度在$\{A\}$中的表示：$\omega_A$
\end{itemize}

\textbf{角速度在$\{B\}$中的表示}：
\begin{equation}
\omega_B = R_{BA} \omega_A
\end{equation}

\textbf{反对称矩阵的变换}：
\begin{align}
[\omega_B] &= [R_{BA} \omega_A] \\
&= R_{BA} [\omega_A] R_{BA}^T \quad \text{（使用相似变换）}
\end{align}

\subsection{示例2：空间角速度与体角速度}

\textbf{空间角速度}（spatial angular velocity）$\omega_s$：在固定坐标系中表示。

\textbf{体角速度}（body angular velocity）$\omega_b$：在旋转坐标系中表示。

\textbf{关系}：
\begin{equation}
\omega_s = R \omega_b
\end{equation}

\textbf{反对称矩阵的关系}：
\begin{align}
[\omega_s] &= [R \omega_b] \\
&= R [\omega_b] R^T \quad \text{（使用相似变换）}
\end{align}

\textbf{逆关系}：
\begin{align}
[\omega_b] &= R^T [\omega_s] R \\
&= [R^T \omega_s]
\end{align}

\subsection{示例3：伴随变换中的应用}

在伴随变换中，我们经常遇到：
\begin{equation}
R [p]_{\times} \omega
\end{equation}

\textbf{使用相似变换}：
\begin{align}
R [p]_{\times} \omega &= R [p]_{\times} R^T (R \omega) \\
&= [R p]_{\times} (R \omega) \quad \text{（使用相似变换）} \\
&= (R p) \times (R \omega)
\end{align}

\textbf{另一种形式}：
\begin{equation}
R [p]_{\times} = [R p]_{\times} R
\end{equation}

\subsection{示例4：速度合成中的应用}

在速度合成中，我们经常遇到：
\begin{equation}
R^T [p]_{\times} \omega
\end{equation}

\textbf{使用相似变换}：
\begin{align}
R^T [p]_{\times} \omega &= R^T [p]_{\times} R (R^T \omega) \\
&= [R^T p]_{\times} (R^T \omega) \quad \text{（使用相似变换）} \\
&= (R^T p) \times (R^T \omega)
\end{align}

\textbf{另一种形式}：
\begin{equation}
R^T [p]_{\times} = [R^T p]_{\times} R^T
\end{equation}

\section{几何意义}

\subsection{物理直观}

反对称矩阵$[\omega]$表示绕轴$\omega$的旋转。

相似变换$R [\omega] R^T$表示：
\begin{itemize}
\item 先应用旋转$R^T$（将坐标系转换到新坐标系）
\item 然后应用旋转$[\omega]$（在新坐标系中旋转）
\item 最后应用旋转$R$（转换回原坐标系）
\end{itemize}

结果等价于直接应用旋转$[R \omega]$。

\subsection{坐标变换的角度}

从坐标变换的角度看：
\begin{itemize}
\item $[\omega]$是在原坐标系中表示的旋转
\item $R [\omega] R^T$是在新坐标系中表示的相同旋转
\item $[R \omega]$是新坐标系中旋转轴$\omega$的反对称矩阵
\end{itemize}

两者等价，因为旋转的物理效果相同。

\section{重要推论}

\subsection{推论1：旋转保持叉积}

从相似变换可以推导出旋转保持叉积：
\begin{align}
R(a \times b) &= R [a] b \\
&= R [a] R^T (R b) \\
&= [R a] (R b) \quad \text{（使用相似变换）} \\
&= (R a) \times (R b)
\end{align}

\subsection{推论2：反对称矩阵的相似性}

所有反对称矩阵在相似变换下保持反对称性：
\begin{equation}
(R [\omega] R^T)^T = R [\omega]^T R^T = -R [\omega] R^T
\end{equation}

因此$R [\omega] R^T$也是反对称矩阵。

\subsection{推论3：特征值的保持}

相似变换保持特征值，因此：
\begin{equation}
\text{eig}(R [\omega] R^T) = \text{eig}([\omega])
\end{equation}

\section{总结}

\subsection{主要公式}

\begin{enumerate}
\item \textbf{基本相似变换}：
\begin{equation}
R [\omega] R^T = [R \omega]
\end{equation}

\item \textbf{逆相似变换}：
\begin{equation}
R^T [\omega] R = [R^T \omega]
\end{equation}

\item \textbf{左乘形式}：
\begin{equation}
R [\omega] = [R \omega] R
\end{equation}

\item \textbf{转置形式}：
\begin{equation}
R^T [\omega] = [R^T \omega] R^T
\end{equation}
\end{enumerate}

\subsection{关键要点}

\begin{itemize}
\item 相似变换将反对称矩阵从一个坐标系转换到另一个坐标系
\item 旋转矩阵保持叉积的性质是相似变换的基础
\item 相似变换在机器人学中用于坐标变换和速度合成
\item 所有形式都是等价的，可以根据需要选择合适的形式
\end{itemize}

\subsection{应用场景}

\begin{itemize}
\item 角速度的坐标变换
\item 空间角速度与体角速度的转换
\item 伴随变换的推导
\item 速度合成公式的简化
\item 动力学方程中的坐标变换
\end{itemize}

\end{document}

